\chapter*{Введение}
\addcontentsline{toc}{chapter}{Введение}
\label{ch:intro}

Видеоплатформы — один из самых нагруженных классов веб-приложений: они
объединяют загрузку и хранение больших файлов, потоковую отдачу видео,
управление пользователями и интерактивный пользовательский интерфейс.
На примере такого приложения удобно показать полный современный стек
веб-разработки: серверный \textbf{API} на \textbf{Python}, клиентское
одностраничное приложение на \textbf{React}, заранее спроектированный
в \textbf{Figma} интерфейс и развёртывание готового продукта на
реальном сервере.

Расчётно-графическая работа посвящена созданию тестового видеохостинга.
В проект заложены основные функции хостинга роликов: регистрация и
вход в аккаунт, загрузка видео, лента опубликованных роликов, страница
просмотра с потоковой отдачей файла и поддержкой перемотки, личный
профиль с управлением собственным контентом.

Дизайн-система интерфейса заранее проработана в Figma: поля ввода с
состояниями, карточки видео, кнопки, шапка приложения, экраны входа и
регистрации. Дизайн-токены (цвета, типографика, отступы) вынесены в
один CSS-файл, что позволяет легко перекрасить интерфейс.

Цель работы — спроектировать и реализовать полнофункциональный
видеохостинг на современном стеке веб-технологий, обеспечить
взаимодействие серверной и клиентской частей через REST-API и
развернуть готовый продукт на сервере с доступом по защищённому
протоколу HTTPS.

Задачи:
\begin{enumerate}
  \item Спроектировать в Figma макет интерфейса: цветовую палитру,
        типографику, состав экранов и компонентов.
  \item Разработать серверную часть на Django: модель данных,
        авторизацию по паролю и гостевые аккаунты с выдачей
        JWT-токенов, загрузку и потоковую отдачу видео, REST-API.
  \item Разработать клиентское приложение на React с маршрутизацией,
        интеграцией с API и реализацией компонентов дизайн-системы.
  \item Покрыть тестами серверную и клиентскую части.
  \item Развернуть платформу на сервере за обратным прокси Caddy,
        оформить документацию API в Git.
\end{enumerate}

\textbf{Инструменты и средства:}
\begin{itemize}
  \item \texttt{Python~3} с фреймворками \texttt{Django} и
        \texttt{Django REST Framework} — серверная часть и REST-API;
  \item \texttt{Gunicorn} — WSGI-сервер;
  \item \texttt{PostgreSQL} — база данных продакшена,
        \texttt{SQLite} — база данных локальной разработки;
  \item \texttt{Node.js}, \texttt{React}, \texttt{React Router},
        \texttt{Vite} — клиентское одностраничное приложение;
  \item \texttt{Figma} — источник макета и дизайн-системы интерфейса;
  \item \texttt{Docker}, \texttt{docker-compose}, обратный прокси
        \texttt{Caddy} с автоматическим TLS Let's~Encrypt —
        развёртывание на собственном сервере;
  \item \texttt{pytest} и \texttt{Vitest} — тестирование серверной и
        клиентской частей;
  \item \texttt{Git} — контроль версий и документация проекта.
\end{itemize}

\endinput
